\setcounter{chapter}{1}
\chaptermark{费马原理与变折射率光学}
\pdfbookmark[0]{第一章 费马原理与变折射率光学}{chapter1}
\problem{1-1}{
可见光谱区在真空中的波长范围一般认定为 $380$--$760\,\mathrm{nm}$，现针对其紫端 $\lambda_1=380\,\mathrm{nm}$、中部 $\lambda_2=550\,\mathrm{nm}$ 和红端 $\lambda_3=760\,\mathrm{nm}$，
\begin{enumerate}
\item 算出这三种光波的时间频率 $f_1$、$f_2$ 和 $f_3$。
\item 当它们传播于折射率为 $1.33$ 的水中时，光波长 $\lambda'_1$、$\lambda'_2$ 和 $\lambda'_3$ 各为多少？其光速 $v'$ 为多少？（忽略色散）
\item 当它们传播于折射率为 $1.58$ 的玻璃中时，光波长 $\lambda''_1$、$\lambda''_2$ 和 $\lambda''_3$ 各为多少？其光速 $v''$ 为多少？（忽略色散）
\end{enumerate}
}
\problem{1-2}{
试从具体考察光程出发，论证：
\begin{enumerate}
\item 反射定律给出的反射光束方向满足等光程性。
\item 折射定律给出的折射光束方向满足等光程性。
\end{enumerate}
提示：参考正文图 1.12。
}
\problem{1-3}{
巨蟹星座中心有一颗脉冲星，其辐射的光频讯号和射频讯号到达地球有 $1.27\,\mathrm{s}$ 的时差，且光波快于射电波。
\begin{enumerate}
\item 求这两种电磁波从脉冲星到地球的光程差 $\Delta L$。
\item 试估算这两种电磁波传播于宇宙空间中的折射率之差 $\Delta n$ 的数量级，以及相应的速度差 $\Delta u/c$ 之数量级。已知这颗脉冲星与地球之间距离为 $D\approx6300$ 光年 $\approx6\times10^{16}\,\mathrm{km}$。
\end{enumerate}
}
\problem{1-4}{
如图所示，一透明平板其厚度为 $h$、折射率为 $n$。一光线从空气射向平板，经其上、下表面反射而分成两条光线，再经透镜而相交于 $P$ 点，试导出光程差公式
\[
L(QABCP)-L(QAP)=2nh\cos i.
\]
这里，$i$ 为折射角。
\tikzfig{fig_1_4}{1-4}
}
\problem{1-5}{
试从费马原理出发导出球面镜反射傍轴成像公式，
\[
\frac{1}{s'}+\frac{1}{s}=-\frac{2}{r}.
\]
（要求作出示意图。）
}
\problem{1-6}{
如图所示，宽度为 $d$ 的一玻璃块其折射率随高度 $y$ 而变化，
\[
n(y)=\frac{n_0}{1-y/r_0},
\]
其中 $n_0=1.2$，$r_0=13\,\mathrm{cm}$。一光线沿 $x$ 轴射向原点而进入这玻璃块，最终从 $A$ 点出射，其倾角为 $\alpha=30^\circ$。试求：
\begin{enumerate}
\item 这条光线在玻璃块中的径迹；
\item 玻璃块在出射点 $A$ 处的折射率 $n_A$；
\item 玻璃块的宽度 $d$。
\end{enumerate}
\tikzfig{fig_1_6}{1-6}
}
\problem{1-7}{
如图所示，一点源 $S$ 发出球面波，经透镜聚焦于 $S'$ 点。一光阑插入此光场，光阑上开有两个小孔 $O$ 和 $Q$，它俩作为次波源发出次波而到达像面。距离 $\overline{OS'}$、$\overline{QS'}$ 分别表示为 $z_0$ 和 $z$；距离 $\overline{OP}$、$\overline{QP}$ 分别表示为 $r_0$ 和 $r$。设光程差 $(z-z_0)=\lambda/6$，光程差 $(r_0-r)=10\frac14\lambda$。问：
\begin{enumerate}
\item 两个次波源 $O$ 与 $Q$ 之间是否有相位差？如是，试求之；
\item 到达 $P$ 点的两个次波扰动之间是否有相位差？如是，试求之。
\end{enumerate}
\tikzfig{fig_1_7}{1-7}
}
\problem{1-8}{
如图所示，由变折射率材料制成的微透镜被用以聚焦平行光束，其折射率分布 $n(r)$ 关于 $z$ 轴对称。现要求其焦距为 $f$，且 $f$ 值远大于微透镜之孔径 $a$，而 $a$ 又远大于微透镜之厚度 $d$。
\begin{enumerate}
\item 试定性分析此折射率变化函数 $n(r)$ 随 $r$ 增加是降低还是提高？
\item 试定量导出 $n(r)$ 函数；设轴上折射率值为 $n_0$。
\end{enumerate}
\tikzfig{fig_1_8}{1-8}
}
